\documentclass[12pt]{article}
\usepackage[margin=1in]{geometry}
\usepackage{setspace}
\usepackage{parskip}

\title{Peer Review: Yuxin Yang --- Second Draft}
\author{Alejandro De La Torre}
\date{March 2026}

\begin{document}

\maketitle

\singlespacing


\textbf{Research Question and Motivation}

To my understanding, the paper studies how rank-based score conversion in China’s National College Entrance Examination reform may distort students’ elective subject choices. As I understand it, the key idea is that once scores are converted by percentile rank within a subject cohort, the payoff to choosing a subject depends not only on a student’s own ability, but also on who else chooses that subject. That creates a strategic sorting problem: students may avoid subjects that attract stronger peers, which can then reinforce imbalance across subjects. I found the motivation compelling because the project connects institutional design to both efficiency and fairness in a very concrete educational setting. More specifically, the paper asks whether flexible subject choice under rank-based conversion generates endogenous imbalance, and whether alternative rules can mitigate those distortions.



\textbf{Suggestions for Improvement}

I think this is a very interesting and ambitious theory paper, and I like that the mechanism is clear at a high level: students are not only choosing subjects they are good at, but also choosing which competitive environment to enter. To me, the strongest part of the draft is that it identifies a real institutional feature---rank-based score conversion---and shows why it can create a fixed-point problem in subject choice. At the same time, I found myself wanting a bit more intuition in the presentation, especially before the formal model begins.

My first suggestion would be to spend a little more time translating the model into plain language before moving into the notation. I think the introduction does a good job motivating the broad issue, but once the model starts, the reader is asked to absorb a lot of technical structure very quickly: latent abilities, noisy signals, realized performance, endogenous cohort CDFs, and equilibrium as a fixed point. I think the paper would become much easier to follow if, before Section 3, you briefly walked the reader through a simple example of the strategic logic. For instance, you might describe a student who is moderately strong in two subjects and explain why that student may avoid the subject expected to attract stronger competitors, even if it is one of her better subjects in absolute terms. I think an example like that would help the reader keep the mechanism in mind once the formal setup begins.

Second, I found myself wondering a bit more about the exact institutional mapping between the real NCEE system and the model. The draft is clearly inspired by the reform, but at the moment the model feels intentionally stylized. I think that is perfectly fine, but I would suggest being very explicit about which parts are meant to capture the real system closely and which parts are abstractions for tractability. For example, the paper uses four subjects and bundles of size two. If that is just a simplified representation, I think it would help to say so directly and explain why that simplification still preserves the main strategic force you care about.

Third, I think the paper could say a bit more about what exactly counts as ``fairness'' and ``welfare'' in the later institutional comparison. Right now those terms are motivating the project, but I was not yet sure how they would be measured once the comparison is done. I think this matters because different notions of fairness may point in different directions. One could care about expected score outcomes, risk exposure, misallocation relative to comparative advantage, dispersion across students, or protection for students in thinner subjects. I would find it helpful if the paper previewed which welfare and fairness objects it ultimately plans to evaluate.

Fourth, I would encourage clarifying the role of beliefs in the equilibrium story. As I understand it, students choose based on expected cohort composition, and equilibrium requires those expectations to be consistent with realized subject shares. I think that logic is central and quite interesting, but I would have appreciated a bit more verbal explanation of what students are assumed to know when they choose. Are they responding to the full population distribution and the scoring rule? Are they assumed to know the distribution of signals in the population? I am not asking for more technical detail here so much as one or two sentences of intuition about the informational environment.

I also wondered whether the paper could say a little more about what results you expect from the model before the conclusion. Right now the conclusion says that the draft builds the theoretical foundation for later comparison, which makes sense, but I think the reader would benefit from a slightly sharper preview of the kinds of comparative statics or institutional tradeoffs the model is likely to deliver. For example, do you expect more flexibility to improve matching but worsen strategic crowding? Do safeguards for low-participation subjects reduce distortion but introduce other inefficiencies? Even a brief roadmap of the likely insights would help the reader see where the theory is going.

Finally, I think a bit of light exposition work would make the draft feel substantially stronger. I found the introduction conceptually strong, but a few sentences were hard for me to parse on a first read. I think small revisions for clarity---especially in the opening paragraph and in the transition from motivation to model---would help the reader more quickly see the value of the project. Overall, though, I think the central idea is strong, and I especially like that the paper is trying to connect a formal strategic model to a concrete policy design problem.


\end{document}
